% ch07-meetings-realtime
\section{الجلسات المرئيّة والزمن الحقيقيّ}

في الفصول السابقة رأينا كيف يُقسَّم النظام إلى طبقات، وكيف يمرّ كل طلب عبر
\en{CQRS}، وكيف تُنشِر العملياتُ أحداثًا دوميّة، وكيف نعزل العالَم الخارجيّ خلف
منافذ ومحوّلات (\en{ports/adapters}). في هذا الفصل نطبّق كل ذلك على ميزتين
«حيّتين» بطبيعتهما: \textbf{مكالمة الفيديو} بين الطالب والمستشار، و\textbf{الدفع
اللحظيّ} للإشعارات والمحادثة عبر \en{SignalR}. الميزتان مختلفتان في التقنية لكنهما
تتشاركان الفكرة نفسها: النظام لا «يستضيف» الوسائط بنفسه، بل ينسّق بين المستخدم
وخدمة سحابية متخصّصة، ويحتفظ هو فقط بالحقائق المهمّة (مَن حضر؟ متى؟ أين التسجيل؟).

\begin{idea}[الفكرة: النظام مُنسِّق لا خادم وسائط]
مكالمة الفيديو تجري بين متصفّحَي المشاركَين وخدمة \en{Azure Communication Services}
(نسمّيها اختصارًا \en{ACS}). خادمُنا \en{.NET} لا يمرّر بايتات الصوت أو الصورة أبدًا؛
كل ما يفعله هو: إنشاء «غرفة»، إصدار رمز دخول قصير العمر لكل طرف، وطلب بدء التسجيل.
هذا يجعل الحِمل على خادمنا خفيفًا جدًّا، ويترك أمور الترميز والشبكة لِمن يُتقِنها.
\end{idea}

\subsection{الغرفة التي بلا خادم: منفذ \en{IMeetingService}}

كل تعامل مع \en{ACS} يمرّ عبر منفذ واحد اسمه \code{IMeetingService}. في الإنتاج
يُختار المحوّل \code{AzureCommunicationMeetingService} تلقائيًّا حين يكون مفتاح
الإعداد \code{Acs:ConnectionString} مضبوطًا؛ وفي التطوير — حيث لا نملك اشتراك
\en{ACS} — يُحقَن بديل بسيط اسمه \code{StubMeetingService} يُرجِع قيمًا وهميّة حتى
تعمل الشاشات بلا اتصال حقيقيّ. هذا هو نفس نمط المنافذ الذي رأيناه من قبل: تبديل
المزوّد قرارُ \emph{إعداد} لا قرار \emph{كود}.

النقطة المفاجئة أوّل مرّة: \textbf{«الغرفة» ليست كائنًا على الخادم}. الغرفة مجرّد
مُعرِّف مكالمة جماعية في \en{ACS}، أي \en{group-call id} وهو ببساطة \en{GUID}. حين
يقبل المستشار الحجز، نولّد هذا المُعرِّف ونحفظه على صفّ الحجز، ولا شيء غير ذلك:

\begin{english}
\begin{lstlisting}[language={[Sharp]C}]
// A booking's "room" is just an ACS group-call id (a GUID).
public Task<MeetingRoom> CreateRoomAsync(Guid bookingId, CancellationToken ct)
    => Task.FromResult(new MeetingRoom(Guid.NewGuid().ToString()));
\end{lstlisting}
\end{english}

هنشوف إيه في السطر ده: الدالة \code{CreateRoomAsync} لا تتّصل بأي خدمة إطلاقًا — إنّها
فقط تُنشئ \en{GUID} جديدًا (\code{Guid.NewGuid}) وتغلّفه في كائن \code{MeetingRoom}
وتُعيده فورًا. لا يوجد نداء شبكة، لا كائن غرفة يُحجَز، لا موارد تُستهلَك مسبقًا.
الغرفة «توجد» بمجرّد أن يعرف الطرفان مُعرِّفها.

\begin{why}
لماذا غرفة بلا خادم؟ لأنّ إنشاء كائنات غرف مسبقًا لكل حجز يعني موارد محجوزة قد لا
تُستخدَم أبدًا (المستشار قد يرفض، أو ينسى الطالب). المُعرِّف الكسول (\en{lazy}) يعني
أنّنا لا ندفع شيئًا حتى لحظة الانضمام الفعليّ. وإن كان المُعرِّف مفقودًا لسبب ما
(سجلّ قديم)، تُنشئ الواجهةُ الخلفية غرفةً جديدة عند الانضمام — بلا تدخّل يدويّ.
\end{why}

\subsection{الانضمام: رمز وصول لكل مشارك ونافذة \en{$\pm$15} دقيقة}

معرفة مُعرِّف الغرفة وحدها لا تكفي للدخول. عند ضغط «انضمّ»، يتحقّق النظام أوّلًا من
\textbf{نافذة الانضمام}: يُسمَح بالدخول فقط ضمن \en{$\pm$15} دقيقة حول موعد الجلسة
المجدوَل (القيمتان \code{JoinOpensBefore} و\code{JoinClosesAfter} كلاهما 15 دقيقة).
قبل ذلك تظهر رسالة «لم تُفتَح بعد»، وبعده «أُغلِقت الغرفة». ثمّ — والأهمّ — يُصدِر
النظام \textbf{رمز وصول مؤقّتًا لكل مستخدم على حدة} من هويّة \en{ACS}:

\begin{english}
\begin{lstlisting}[language={[Sharp]C}]
public async Task<MeetingAccessToken> IssueAccessTokenAsync(
    string roomId, Guid userId, CancellationToken ct)
{
    var result = await _identityClient
        .CreateUserAndTokenAsync(new[] { CommunicationTokenScope.VoIP }, ct)
        .ConfigureAwait(false);

    return new MeetingAccessToken(
        RoomId: roomId,
        Token: result.Value.AccessToken.Token,
        AcsUserId: result.Value.User.Id,
        ExpiresAt: result.Value.AccessToken.ExpiresOn);
}
\end{lstlisting}
\end{english}

هنشوف إيه هنا سطرًا-بسطر:
\begin{itemize}
\item \code{CreateUserAndTokenAsync} تُنشئ هويّة \en{ACS} مؤقّتة \emph{وتُصدِر لها رمزًا}
  في نداء واحد. النطاق المطلوب هو \code{CommunicationTokenScope.VoIP} فقط — أي «صوت
  وفيديو»، لا صلاحيّة أوسع.
\item الرمز الناتج (\code{AccessToken.Token}) هو ما تأخذه الواجهةُ لتدخل به المكالمة،
  ومعه \code{AcsUserId} الذي يميّز المشارك داخل \en{ACS}.
\item \code{ExpiresOn} تحمل لحظة انتهاء الصلاحيّة: الرمز قصير العمر (مدّته نحو ساعتين)،
  وهي أكثر من كافية لجلسة أقصاها 60 دقيقة.
\end{itemize}

\begin{secnote}
\textbf{رمز الفيديو = بيانات اعتماد.} أي شخص يحمل الرمز يستطيع الدخول للمكالمة، ولهذا
نُصدِر رمزًا \emph{لكل مشارك} مقصورًا على \en{VoIP} وينتهي سريعًا. النتيجة: لو تسرّب
رمزٌ، فهو عديم الفائدة بعد انتهاء الجلسة، ولا يمكن «إعادة تشغيله» في مكالمة أخرى.
لاحظ أيضًا أنّنا لا نسلّم المُعرِّف والرمز إلّا للطرفين المصرَّح لهما عبر نقطة
الانضمام المحميّة — لا عبر أي رابط عامّ.
\end{secnote}

\begin{qa}
\textbf{سؤال متوقّع:} «لماذا لا تُصدِرون رمزًا واحدًا للغرفة وتشاركونه؟» لأنّ الرمز
المشترَك لا يمكن إبطاله لطرفٍ دون الآخر، ولا يميّز مَن دخل. الرمز لكل مستخدم يربط
الدخول بهويّة، ويجعل تتبّع الحضور دقيقًا، ويقلّص أثر أي تسريب إلى مستخدم واحد ولمدّة
قصيرة فقط.
\end{qa}

\subsection{التسجيل والحضور و«الغياب»}

بمجرّد أن تتّصل المكالمة فعليًّا، تطلب الواجهةُ بدءَ التسجيل. التسجيل في \en{ACS}
يتمّ عبر خدمة \en{Call Automation}، ونادينا لها بسيط:

\begin{english}
\begin{lstlisting}[language={[Sharp]C}]
public async Task<string> StartRecordingAsync(string serverCallId, CancellationToken ct)
{
    var result = await _callAutomationClient.GetCallRecording()
        .StartAsync(new StartRecordingOptions(new ServerCallLocator(serverCallId)), ct)
        .ConfigureAwait(false);
    return result.Value.RecordingId;   // handle to the running recording
}
\end{lstlisting}
\end{english}

هنشوف إيه: نمرّر \code{serverCallId} — وهو مُعرِّف المكالمة الجارية كما يراه الخادم —
إلى \code{StartAsync}، فيبدأ \en{ACS} بتسجيل الجلسة إلى تخزين \en{Blob}، ويُعيد لنا
\code{RecordingId} وهو «مقبض» للتسجيل الجاري نحتفظ به. بعد انتهاء الجلسة، نحفظ صفّ
\code{SessionRecording} يحمل مسار التخزين وحجم الملفّ، ونستطيع لاحقًا تنزيل التسجيل
بثًّا عبر \code{DownloadRecordingAsync}، فيظهر للطالب زرّ تنزيل \en{MP4}.

أمّا \textbf{الحضور}، فهو حقائق نحرص على تسجيلها على صفّ الحجز نفسه:
\begin{itemize}
\item \code{StudentJoinedAt} و\code{ConsultantJoinedAt}: لحظة أوّل دخول لكل طرف
  (متطلَّب \en{FR-217}).
\item \code{RecordingStartedAt}: لحظة بدء التسجيل (\en{PB-006}).
\item \code{NoShowMarkedAt}: تُضبَط حين تكتشف مهمّةٌ خلفيّة أنّ طرفًا لم يحضر.
\end{itemize}

\begin{keyp}[نقطة مهمة: مَن يكتشف الغياب؟]
مهمّة خلفيّة دوريّة اسمها \code{MeetingNoShowSweepJob} تمرّ على الحجوزات المنتهية،
فتقارن أوقات الدخول: إن لم يحضر طرفٌ ضِمن النافذة، تُعلِّم الحجز بـ\en{no-show}
(تضبط \code{NoShowMarkedAt}) وتُطلِق منطق الاسترداد المناسب. أي أنّ «الغياب» ليس
حُكمًا يدويًّا بل نتيجة قياس آليّ على أوقات دخول مسجَّلة بدقّة.
\end{keyp}

\subsection{التدفّق الكامل في 12 خطوة}

الجدول التالي يجمع كل ما سبق في رحلة واحدة، من لحظة قبول الحجز حتى تنزيل ملفّ
\en{MP4}. المراحل ملوّنة ذهنيًّا: \textbf{القبول} (خطوات 1–3)، \textbf{التحقّق
والرمز} (4–7)، \textbf{الواجهة و\en{SDK}} (8–10)، ثمّ \textbf{التسجيل والحفظ}
(11–12).

\begin{center}
\begin{tabular}{C{0.9cm} L{13.0cm}}
\toprule
\textbf{\color{spPrimary}\#} & \textbf{\color{spPrimary}الخطوة} \\
\midrule
1 & المستشار يقبل الحجز، فيعمل الأمر \code{AcceptBookingCommand}. \\
2 & المعالِج ينادي \code{IMeetingService.CreateRoomAsync} فيتولّد \code{roomId} من نوع \en{GUID}. \\
3 & يُحفَظ \code{MeetingRoomId} في \en{SQL}، ويُنشَر \code{BookingConfirmedEvent} فيصل إشعار للطالب. \\
4 & ضمن نافذة \en{$\pm$15} دقيقة من الموعد، يضغط الطالب أو المستشار «انضمّ». \\
5 & الأمر \code{RecordMeetingJoinCommand} يتحقّق من النافذة، ثمّ ينادي \code{IssueAccessTokenAsync}. \\
6 & \en{ACS} عبر \code{CreateUserAndTokenAsync} بنطاق \en{VoIP} يُصدِر رمزًا قصير العمر (نحو ساعتين). \\
7 & الواجهةُ تستقبل \code{roomId} و\code{accessToken} و\code{acsUserId}. \\
8 & \en{ACS Calling SDK}: إنشاء \code{CallClient} ثمّ \code{createCallAgent} ثمّ \code{agent.join} بـ\code{groupId}. \\
9 & طلب أذونات الأجهزة، تعداد الكاميرات، إنشاء \code{LocalVideoStream} وعرضه. \\
10 & عند \code{Connected}، تنادي الواجهةُ نقطةَ بدء التسجيل مرسِلةً \code{serverCallId}. \\
11 & \code{StartRecordingAsync} يبدأ التسجيل إلى \en{Blob} ويُعيد \code{RecordingId}. \\
12 & بعد الجلسة: يُحفَظ صفّ \code{SessionRecording}، ويستطيع الطالب تنزيل \en{MP4}. \\
\bottomrule
\end{tabular}
\end{center}

\begin{note}
لاحظ أنّ التسجيل يبدأ من \emph{الواجهة} عند لحظة الاتصال، لا من الخادم مسبقًا. هذا
مقصود: لا معنى لتسجيل غرفة فارغة، والبدء عند \code{Connected} يضمن أنّ هناك مكالمة
حيّة فعلًا. والنداء مُصمَّم ليكون آمنًا للتكرار (\en{idempotent})، فلو أطلقه الطرفان
معًا لا يبدأ تسجيلان.
\end{note}

\subsection{الواجهة: تكامل \en{ACS Calling SDK} في \en{Meeting.tsx}}

الآن نرى الجانب الآخر من نقطة الانضمام: ملفّ \code{Meeting.tsx} في الواجهة. بعد أن
يوافق المستخدم على إشعار التسجيل (بوّابة الموافقة \en{CONSENT-01})، يبدأ الكود
الحقيقيّ التالي (مختصَر):

\begin{english}
\begin{lstlisting}[language=Java]
const join = await meetingsApi.join(bookingId!);
// When ACS isn't configured the server issues a stub token.
if (join.provider === "Stub") {
  throw new Error(t("bookings:meeting.errorNotConfigured"));
}

const callClient = new CallClient();
const credential = new AzureCommunicationTokenCredential(join.accessToken);
const agent = await callClient.createCallAgent(credential, { displayName });
// ... ask device permissions, pick a camera, build localStream ...

const call = agent.join(
  { groupId: join.roomId },
  { videoOptions: { localVideoStreams: localStream ? [localStream] : [] },
    audioOptions: { muted: false } },
);

// Start the recording once the call is actually connected (idempotent).
call.on("stateChanged", () => {
  if (call.state === "Connected") void ensureRecording();
});
\end{lstlisting}
\end{english}

هنشوف إيه في هذا المقتطف:
\begin{itemize}
\item \code{meetingsApi.join} يطلب من الخادم بيانات الدخول؛ لو كان المزوّد
  \code{Stub} (بيئة تطوير) نرمي رسالة صريحة «الفيديو غير مُهيَّأ» بدل خطأ غامض من
  الـ\en{SDK}.
\item ننشئ \code{CallClient} ونغلّف الرمز في \code{AzureCommunicationTokenCredential}،
  ثمّ نصنع «وكيلًا» (\code{createCallAgent}) هو مَن يجري المكالمات باسمنا.
\item الانضمام هو \code{agent.join} مع \code{groupId: join.roomId} — هنا يظهر بوضوح
  أنّ الغرفة مجرّد مُعرِّف مجموعة، ونمرّر خيارات الفيديو والصوت.
\item نستمع لتغيّر الحالة، وعند \code{Connected} فقط نستدعي \code{ensureRecording}
  التي تجلب \code{serverCallId} وتطلب بدء التسجيل من الخادم.
\end{itemize}

\begin{secnote}
لاحظ بوّابة \en{CONSENT-01}: لا ينضمّ المستخدم — ومن ثمّ لا يبدأ أي تسجيل — قبل أن
يوافق صراحةً على أنّ الجلسة مسجَّلة. المؤثّر (\en{effect}) الذي يشغّل الانضمام لا
يعمل إلّا بعد قلب علامة الموافقة إلى «نعم»، فلا شيء يدخل الغرفة قبل الموافقة. هذه
حماية للخصوصيّة مبنيّة في التصميم لا مضافة لاحقًا.
\end{secnote}

\subsection{الزمن الحقيقيّ: \en{SignalR} وثلاثة \en{hubs}}

الجزء الثاني من الفصل يخصّ الدفع اللحظيّ للأحداث النصّية: رسالة محادثة جديدة، إشعار
حجز، تفاعل مجتمعيّ. هنا نستخدم \en{SignalR} عبر ثلاثة \en{hubs}، كلّها ترث قاعدةً
واحدة مُوثَّقة (\code{AuthenticatedHub} المعلَّمة بـ\code{[Authorize]})، فلا يتّصل
بها إلّا مستخدم يحمل \en{JWT} صالحًا. الجدول يلخّص مَن يفعل ماذا:

\begin{center}
\begin{tabular}{L{3.0cm} L{6.6cm} L{4.0cm}}
\toprule
\textbf{\color{spPrimary}الـ\en{hub}} & \textbf{\color{spPrimary}الأحداث المدفوعة} & \textbf{\color{spPrimary}من يستخدمه} \\
\midrule
\code{ChatHub} & رسالة جديدة، مؤشّر الكتابة، الحضور (\en{online/offline}) & محادثات المجتمع الخاصّة \\
\code{NotificationHub} & تأكيد حجز، دفعة معالَجة، تذكير موعد & كل الأدوار \\
\code{CommunityHub} & ردود وتفاعلات جديدة، انضمام للأقسام والمواضيع & موجز المجتمع \\
\bottomrule
\end{tabular}
\end{center}

كيف يصل الحدث إلى مستخدم بعينه؟ الفكرة أنّ كل اتّصال جديد يُضاف إلى «مجموعة» خاصّة
بصاحبه. في \code{NotificationHub} مثلًا، عند الاتّصال نضيف الاتّصال إلى مجموعة اسمها
\code{user:<userId>}:

\begin{english}
\begin{lstlisting}[language={[Sharp]C}]
public override async Task OnConnectedAsync()
{
    // Await the group-add: a notification fired right after connect must not
    // race past the join and miss the recipient.
    if (Context.UserIdentifier is not null)
    {
        await Groups.AddToGroupAsync(
            Context.ConnectionId, $"user:{Context.UserIdentifier}");
    }
    await base.OnConnectedAsync();
}
\end{lstlisting}
\end{english}

هنشوف إيه: \code{Context.UserIdentifier} هو مُعرِّف المستخدم المستخرَج من الـ\en{JWT}.
نضيف اتّصاله الحاليّ إلى مجموعة \code{user:<id>}، \emph{وننتظر} إتمام الإضافة
(\code{await}) لا نطلقها ونمضي؛ لأنّ إشعارًا يُدفَع بعد الاتّصال مباشرةً قد «يسبق»
انضمام المجموعة فيضيع على صاحبه ثمّ يظهر لاحقًا كأنّه مكرَّر.

بعد ذلك، الدفعُ نفسه سطرٌ واحد من أي خدمة تحمل \code{IHubContext}. هذا مثال حقيقيّ
من موزّع الإشعارات، يدفع للمجموعة الخاصّة بالمستخدم المستهدَف:

\begin{english}
\begin{lstlisting}[language={[Sharp]C}]
await hub.Clients.Group($"user:{row.RecipientUserId}")
    .SendAsync("notification", new
    {
        id = row.Id,
        type = row.Type.ToString(),
        row.TitleEn,
        row.TitleAr,
        row.DeepLink,
        createdAt = DateTimeOffset.UtcNow,
    }, ct);
\end{lstlisting}
\end{english}

هنشوف إيه: نأخذ \code{hub} (من نوع \code{IHubContext<NotificationHub>})، نستهدف
مجموعة المستخدم، ونرسل حدثًا باسم \code{notification} يحمل حمولة \en{JSON} بسيطة.
على الواجهة، خطّاف \en{React} يستمع لهذا الاسم فيحدّث الحالة فورًا — فيظهر الإشعار
دون أن يُحدِّث المستخدم الصفحة. لاحظ أنّنا نرسل العنوان بالعربية والإنجليزية معًا
(\code{TitleAr} و\code{TitleEn})، فتختار الواجهةُ لغة العرض.

\begin{why}
لماذا مجموعة \code{user:<id>} بدل استهداف اتّصال بعينه؟ لأنّ للمستخدم الواحد قد تكون
عدّة تبويبات أو أجهزة متّصلة في آنٍ. الدفع للمجموعة يصل إلى \emph{كل} جلساته الحيّة
مرّة واحدة، وهو المكافئ العمليّ لِنمط \code{Clients.User(id).SendAsync}. وحين نحتاج
دفعًا لطرفَي محادثة أو لمشتركي موضوع، نستخدم مجموعات أخرى مثل
\code{conversation:<id>} و\code{forum-post:<id>}.
\end{why}

\begin{secnote}
الاشتراك في المجموعات ليس مفتوحًا: قبل أن ينضمّ عميلٌ لمجموعة محادثة أو موضوع، يتحقّق
الـ\en{hub} أنّ صاحب الـ\en{JWT} طرفٌ فعليّ فيها (بوّابة \en{SEC-02})، وإلّا يرمي
\en{HubException}. بغير هذا التحقّق كان أي مستخدم مصادَق يستطيع الاشتراك في
\code{conversation:<id>} والتنصّت على رسائل خاصّة — لذا الاشتراك نفسه محميٌّ لا الدفع
فقط.
\end{secnote}

\begin{summary}[الخلاصة في سطور]
\begin{itemize}
\item مكالمة الفيديو تمرّ عبر منفذ \code{IMeetingService} (محوّل
  \code{AzureCommunicationMeetingService} في الإنتاج، \code{StubMeetingService} في
  التطوير)، ويُختار بالإعداد \code{Acs:ConnectionString}.
\item «الغرفة» = \en{ACS group-call id} أي \en{GUID}، بلا كائن خادم؛ تُنشأ كسولًا
  عند القبول.
\item عند الانضمام: نافذة \en{$\pm$15} دقيقة، ورمز وصول \en{VoIP} قصير العمر (نحو ساعتين)
  \emph{لكل مشارك} عبر \code{CreateUserAndTokenAsync} — لأنّ رمز الفيديو بيانات
  اعتماد.
\item التسجيل عبر \en{Call Automation} (\code{StartRecordingAsync} \en{$\rightarrow$}
  \code{RecordingId}) يُحفَظ كـ\code{SessionRecording} + \en{Blob}؛ والحضور
  (\code{StudentJoinedAt}/\code{ConsultantJoinedAt}) يقوده \code{MeetingNoShowSweepJob}
  لاكتشاف \en{no-show}.
\item الواجهة \code{Meeting.tsx} تنضمّ بـ\code{groupId}، تعرض الكاميرا، وتبدأ التسجيل
  عند \code{Connected} خلف بوّابة موافقة \en{CONSENT-01}.
\item الزمن الحقيقيّ عبر ثلاثة \en{hubs} (\code{ChatHub}/\code{NotificationHub}/%
  \code{CommunityHub}) مصادَقة بـ\en{JWT}؛ الدفع سطرٌ واحد عبر \code{IHubContext}
  إلى مجموعة \code{user:<id>}، والاشتراك في المجموعات محميّ (\en{SEC-02}).
\end{itemize}
\end{summary}
